% !TEX root = ../thesis.tex

\chapter{Úvod}\label{ch:introduction}

V súčasnej dobe sa náhodnosť stáva jedným z kľúčových prvkov v oblasti informatiky, kryptografie, simulácií a modelovania fyzikálnych procesov. Náhodné čísla predstavujú nevyhnutnú súčasť mnohých moderných technológií a systémov, v ktorých sa využívajú na zabezpečenie bezpečnej komunikácie, autentifikácie používateľov či na simulácie komplexných prírodných alebo technických procesov. Ich úloha je zásadná najmä v kryptografii, kde kvalita a nepredvídateľnosť náhodných čísel priamo určujú úroveň bezpečnosti šifrovacích algoritmov. Bez spoľahlivého zdroja náhodnosti by bolo prakticky nemožné garantovať dôvernosť, integritu ani overiteľnosť prenášaných dát. Rovnako aj v oblasti vedeckých výpočtov a simulácií náhodné čísla umožňujú modelovať správanie systémov, ktoré sú ovplyvnené náhodnými alebo štatistickými procesmi, čím poskytujú nástroj pre realistické experimentálne overovanie hypotéz.

Napriek ich zásadnej dôležitosti sa väčšina dnešných aplikácií spolieha na tzv. pseudonáhodné generátory, ktoré vytvárajú sekvencie čísel prostredníctvom deterministických matematických algoritmov. Tieto generátory síce dokážu produkovať čísla, ktoré sa na prvý pohľad javia ako náhodné, no ich vnútorná štruktúra je úplne predvídateľná, ak poznáme počiatočné podmienky alebo algoritmus generovania. To znamená, že ide iba o „zdanlivú“ náhodnosť, ktorá nevychádza z fyzikálne nepredvídateľných procesov, ale z deterministického výpočtu. Takéto obmedzenia môžu viesť k opakovaniu vzorov, nižšej entropii a tým k zníženiu bezpečnosti systémov, ktoré sú na náhodnosť citlivé – napríklad pri generovaní šifrovacích kľúčov alebo autentifikačných tokenov. Nedokonalosť pseudonáhodných generátorov preto predstavuje reálny problém, ktorý je potrebné riešiť, najmä v prostredí, kde je absolútna náhodnosť nevyhnutnou podmienkou bezpečnosti a spoľahlivosti.

Kvantová mechanika prináša nový pohľad na pojem náhodnosti. Na rozdiel od klasickej fyziky, ktorá je deterministická, kvantové javy sú prirodzene pravdepodobnostné. Tento fakt umožňuje konštruovať zariadenia, ktoré generujú skutočne náhodné čísla na základe merania kvantových stavov. Medzi najvýznamnejšie kvantové princípy využiteľné pre generovanie náhodnosti patria superpozícia a prepletenie (entanglement). Najmä prepletené stavy – známe ako Bellove páry – predstavujú základ pre teoreticky dokonalý zdroj náhodnosti. Meranie takýchto stavov podľa princípov kvantovej mechaniky produkuje výsledky, ktoré nie je možné predpovedať ani replikovať, a tým vytvárajú ideálny zdroj skutočne náhodných čísel.

V tejto práci vychádzam zo tzv. Stockholmskej interpretácie kvantovej mechaniky, ktorá považuje výsledok merania za konečný a fyzikálne zmysluplný jav.\cite{hanson1959copenhagen} Práca sa preto sústreďuje výhradne na tú časť kvantového procesu, ktorá sa odohráva po samotnom meraní a ktorá je matematicky opísateľná. Nezaoberám sa otázkami deterministickej alebo nedeterministickej povahy vesmíru ako celku, ani filozofickými implikáciami kolapsu vlnovej funkcie. Cieľom nie je skúmať fundamentálnu podstatu reality, ale analyzovať praktickú využiteľnosť existujúcich kvantových javov ako zdroja náhodnosti, ktorú možno empiricky overiť a technicky implementovať.

Zámerom práce je teda nielen experimentálne potvrdiť, že kvantové generátory dokážu produkovať kvalitnejšie náhodné čísla než klasické algoritmy, ale aj skúmať praktické limity ich použitia. Tieto limity zahŕňajú najmä dostupnosť kvantového hardvéru, finančné náklady spojené s využitím kvantových výpočtových služieb, časovú náročnosť výpočtov a otázku škálovateľnosti pre reálne aplikácie. V neposlednom rade bude analyzovaná aj kvalita získaných údajov pomocou štatistických testov náhodnosti a porovnanie so štandardnými pseudonáhodnými generátormi používanými v bežných systémoch.

Kvantové generátory náhodných čísel predstavujú jednu z mála oblastí kvantového výpočtového prostredia, kde už dnes existuje reálne praktické využitie. Ich potenciál spočíva v schopnosti generovať dáta, ktoré spĺňajú najprísnejšie požiadavky na entropiu a nepredvídateľnosť. Napriek tomu, že teória kvantovej náhodnosti je dobre pochopená, praktické implementácie sú stále vo fáze výskumu a testovania. Táto práca preto prispeje k pochopeniu, ako blízko sú súčasné kvantové technológie k ideálu dokonalého zdroja náhodných čísel, a aké sú ich obmedzenia z pohľadu technickej realizácie a ekonomickej efektivity.




\section*{Formulácia úlohy}

Úlohou tejto bakalárskej práce je navrhnúť, implementovať a experimentálne overiť kvantový generátor náhodných čísel využívajúci merania kvantových stavov s prepletením (Bellove páry). Praktická časť projektu bude realizovaná prostredníctvom cloudovej kvantovej platformy \textit{Amazon Braket}\cite{gonzalez2021cloud}, ktorá umožňuje spúšťať kvantové algoritmy na reálnych kvantových zariadeniach.

Súčasťou riešenia bude vytvorenie softvérového systému založeného na frameworku \textit{Django}~\cite{rubio2017beginning}, ktorý zabezpečí generovanie, zber, ukladanie a analýzu vygenerovaných kvantových údajov. Systém umožní opakované generovanie dávok náhodných čísiel, ich ukladanie do databázy a vykonanie štatistických analýz kvality náhodnosti. Ďalej bude implementovaný modul na porovnanie výsledkov s ideálnym teoretickým modelom, ako aj vizualizačná časť prezentujúca distribúcie hodnôt a odchýlky od rovnomerného rozdelenia.

